\cotefile{93}

\scanpage{76}

\unsure{b'|b)|Si} on s'intéresse à $\unsure{\mathscr E|\mathscr L|\mathcal E}_X$ au lieu de $\underline O_X$, il faut remplacer $\underline\omega_X$ par $\underline\omega_X^*$ (sous-faisceau des bases de $\underline\omega_X$, i.e. des sections sans zéros), qui est un torseur sous $\underline O_X^*$ et $dg-df$ par $\deleted{dg}\mathrel{:} A'(f,g)=\frac{dg}{df}\in\Gamma(X,\underline O_X^{\times})$

b) \uline{$G_m$ opérant sur $\mathbb E^{\unsure{1|\prime}}$ ou $\mathbb E^{1\times}\simeq G_m$} par \deleted{translation}\inserted{homothétie}. Plus généralement, considérons son opération sur un faisceau inversible $L$, soit $L^{\times}$ le sous faisceau des bases \inserted{(torseur sous $\underline O_X^{\times}$)}, on s'intéresse à $L^{\times}/\mathbb C^{\times}$. Or

à \unsure{tte|une} section de $L^{\times}$ est \unsure{associé|associé'|associée} une connexion sur $L$. on trouve ainsi un \unsure{hom.|hom|homom.}

\[
 L^{\times}\,\deleted{\unsure{/}}\xrightarrow{D}\underline{\mathrm{Conn}}_{X/\mathbb C}(L)\qquad(f\longmapsto D_f)
\]

où le deuxième mb est un torseur sous $\underline\omega_{X/\mathbb C}^{\unsure{1}}$, et ceci induit une \deleted{\unsure{bijection}} Iso

\[
 L^{\times}/\mathbb C^{\times}\simeq\underline{\mathrm{Conn}}_{X/\mathbb C}(L)
\]

Si on ignorait la notion de connexion, on introduirait directement $B(f,g)$ $(f,g\in\Gamma L^*)$, qui sera \deleted{\unsure{comme fonct.|conn.}} $D_g-D_f$, par

\[
 B(f,g)=\frac{d\,g/f}{g/f}\in\Gamma\deleted{(}\underline\omega_{X/\mathbb C}
\]

avec des propriétés du type a) b) c) d) e). Dans le cas $L=\underline O_X$, $\underline{\mathrm{Conn}}(L)\simeq\underline\omega^{\unsure{1}}_{X/\mathbb C}$ can$^{\mathrm{t}}$, et il suffit de considérer

\[
 g\longmapsto\frac{dg}{g}\qquad O_X^*\longrightarrow\underline\omega_{X/\mathbb C}
\]

qui induit $\underline O_{X}^{\times}/\mathbb C^{\times}\simeq\underline\omega_{X/\mathbb C}$

% archive sheet number: 75
\scanpage{77}

C'est \unsure{associé|mené} à $B'(f,g)=\frac{dg}{g}-\frac{df}{f}$ .\,.

Si on veut par contre un invariant sur $\unsure{\mathscr E|\mathcal E}_X^*$, on prendra

\[
 g\longrightarrow\frac{g}{dg}\qquad \mathscr E_X\longrightarrow\underline\omega_{X/\mathbb C}^{\otimes(-1)}
\]

on \deleted{trouve} définit

\[
 \deleted{\unsure{B}}C^{!}(f,g)=\deleted{\frac{dg}{g}-\frac{df}{f}}\;g(dg)^{-1}-f(df)^{-1}
\]

satisfaisant à a) b) c) [$C^{?!}(f,g)=C^{?!}(f,ag)$ si $a\in\mathbb C^{\times}$) d) ($C^{?!}(f,g)=C^{?!}(f,h)$ et $g(x)=h(x)$ \uline{en un pt qui n'est pas un zéro} de $g$ ou $h$ $\Rightarrow g=h$ ] et e) \deleted{\unsure{$[g\mapsto C(f,g)$|$[g\mapsto C^!(f,g)]$}} pas tout à fait e, car $g(dg)^{-1}-f(df)^{-1}$ \unsure{p.ex|p.ex.} $g(dg)^{-1}$ \unsure{ne|va} pas être id. nul si $g\in\Gamma\,\mathscr E_X$ ! Donc il faut se borner à

\[
 \mathscr E_X^*=\underline{\mathrm{Ét}}_{\mathbb C\text{-an}}(X,\unsure{G_m^*|G_m^{\times}|\mathbb C^*})
\]

au quel cas on l'envoie dans $\underline\omega_{X/\mathbb C}^{\otimes(\deleted{\unsure{+}}-1)\,*}$ par $g(dg)^{-1}$, ou au choix dans $\underline\omega_{X/\mathbb C}^*$ par $dg/g$, et il y a lieu de considérer $\omega_{X/\mathbb C}^{\times}$ comme un torseur sous $\underline O_X^*$\unsure{,|;|.}

Donc ici on \unsure{poserait|prend}

\[
 C^{\unsure{?|\prime|\natural}}(f,g)=\frac{dg}{g}\mathbin{:}\frac{df}{f}=\frac{dg}{df}\,\frac{g}{f}\in\Gamma(X,\underline O_X^{\times})
\]

Dans tous ces cas, l'espace des valeurs \inserted{$(\mathbb P^1(\mathbb C),\mathbb E^1(\mathbb C),G_m(\mathbb C))$} pour les $f,g$ était un espace \uline{homogène} sous le groupe d'op. envisagé ($\unsure{\mathrm{Gb}(1)|\mathrm{GP}(1)|\mathrm{GL}(1)}$, $G_a$, $G_m$) ce qui semble condition de succès ($\unsure{\mathbb C^{\times}|G_m^{\times}|\mathbb C^*}$ opérant sur $\mathbb C$ : mauvais !) 

c/ Groupe $\mathrm{Aff}(1)$ \deleted{$\unsure{G}$} $=G_m\cdot G_a$ opérant sur $\mathbb E^1$.

On considère alors

% archive sheet number: 76
\scanpage{78}

\diagram{\begin{tikzcd}[ampersand replacement=\&,column sep=small,row sep=small]
\unsure{\mathscr E|\mathcal E}_X \arrow[rr] \arrow[dr] \& \& \underline{\mathrm{Conn}}(\underline\omega_X)\\
\& \underline\omega_{X/\mathbb C}^{*} \arrow[ur] \&
\end{tikzcd}}

par $f\longrightarrow D_{df}$ ($=$ connexion de $\underline\omega_X$ pour laquelle \inserted{la base} $df$ est horizontale)

on trouve

\[ {}_{\mathrm{Aff}(1,\mathbb C)}\backslash\mathscr E_X\simeq\underline{\mathrm{Conn}}(\underline\omega_X) \]

car successivement par a)

\[ {}_{G_a(\mathbb C)}\backslash\mathscr E_X\overset{d}{\simeq}\underline\omega_X^{\times} \]

et par b)

\[ {}_{G_m(\mathbb C)}\backslash\underline\omega_X^{\times}\xrightarrow{\sim}\underline{\mathrm{Conn}}(\underline\omega_X) \]

On trouve donc

\[ D(f,g)=\frac{d\,\frac{dg}{df}}{\frac{dg}{df}}\in\Gamma\,\underline\omega_X \]

Si on a unif$^{\mathrm{te}}$ $t$, cela fait

\[ D(f,g)=\frac{d\,\frac{g'}{f'}}{\frac{g'}{f'}}=\frac{dg'}{g'}-\frac{df'}{f'}=\left(\frac{g''}{g'}-\frac{f''}{f'}\right)dt \]

\deleted{\unsure{par conséquent|Remarquons}}

d) \deleted{$\unsure{\mathscr E_X|\mathcal E_X}$} \uline{Groupe} ``\uline{affin}'' \uline{de} $G_m$

\[ \underset{\substack{\Vert\\ \mathbb Z/2\mathbb Z}}{\mathrm{Aut}(G_m)}\cdot G_m \]

\uline{opérant sur $G_m$.} ?

\[ f\longmapsto\frac{1}{f}\quad\text{et}\quad f\unsure{\mapstochar\relbar|\longmapsto} c f\quad(c\in\mathbb C^{\times}) \]

$f$ dans $\underline{\mathrm{Ét}}(X,\mathbb C^{\times})=\mathscr E_X^{\times}$. On prend

\diagram{\begin{tikzcd}[ampersand replacement=\&,column sep=small,row sep=small]
\mathscr E_X^{\times}\arrow[rr]\arrow[dr,swap,"{f\mapsto\frac{df}{f}}"] \& \& (\underline\omega_X^{\otimes2})^* .\\
\&\underline\omega_X^*\arrow[ur,swap,"{w\mapsto w^2}"]\&
\end{tikzcd}}

% archive sheet number: 77
\scanpage{80}

$\omega$ base de $\underline\omega^{\otimes n}$ ($n\in\mathbb Z$, \unsure{$\exists$|ici} $n\ne0$) $\omega^{1/n}$ base de $\underline\omega$ définie loc. à mult. près $\varepsilon$, $\varepsilon^n=1$ $\int\omega^{1/n}$ bien défini mod transformation $f\longmapsto\varepsilon f+b$ $(\varepsilon^{n}=1)$ donc $c_{\int\omega^{1/n}}$ bien déf$^{\mathrm i}$

Si $t$ uniformisante et $\omega=f(t)dt^n$, on aura $\omega^{1/n}=f^{1/n}dt$ et $\deleted{\unsure{c_{\int\omega^{1/n}}}}\,c_{\int\omega^{1/n}}=c_{\int f^{1/n}dt}\deleted{=c_t+}$

\[ c_n(\omega)=c_t+\left[-\frac2n\frac{f''}{f}+\frac{2n+1}{n^2}\left(\frac{f'}f\right)^2\right]dt^2 \qquad\text{si }n\ne0 \]

\deleted{$\unsure{G_n(\omega)}$}

\[ \begin{aligned}
 G_{m,n}\bigl(\underset{\substack{\Vert\\ f(dt)^m}}{\omega_m},\underset{\substack{\Vert\\g(dt)^{n}}}{\omega_n}\bigr)
 &=c_n(\omega_n)-c_m(\omega_m)=\left[\left(-\frac2n\frac{g''}g+\frac{2n+1}{n^2}\left(\frac{g'}g\right)^2\right)\right.\\
 &\qquad\left.-\left(-\frac2m\frac{f^{\unsure{4|''}}}{f}+\frac{2m+1}{m^2}\left(\frac{g'}g\right)^2\right)\right]dt^2
 \end{aligned} \]

\[ \left\{\begin{aligned}
 G_{m,n}(\omega_m,\omega'_n)+G_{n,\unsure{r|p}}(\omega'_n,\omega''_p)&=G_{m,p}(\omega_m,\omega''_p)\\
 G_{m,n}(\omega_m,\omega_n)&=G_{m,nr}(\omega_m,\omega_n^r)\qquad\text{si }n\ne0,r\ne0\\
 G_{m,n}(\omega_m,ds)&=G_{m,0}(\omega_m,s)\\
 G_{0,0}(t,s)&=G(t,s)
 \end{aligned}\right. \]

\inserted{$G_{m,n}(\omega_m,(ds)^n)=G_{m,0}(\omega_m,s)$}

\[ c'_n(\omega)=n^2\frac{c_n(\omega)}{\deleted{\unsure{12}}}\in \unsure{\mathscr T|\mathcal T}_X^{* n^2} \]

\[ G_{m,n}((dt)^m,(ds)^n)=G_{0,0}(t,s) \]

\[ \deleted{c'_n(\omega)}\qquad m,n\qquad\delta=m\alpha=n\beta\qquad\delta^2=m^2\alpha^2=n^2\beta^2 \]

\[ \begin{aligned}
 \omega_m,&\longmapsto\alpha^2\deleted{\unsure{c}}c'_m(\omega_m)=\alpha^2m^2\frac{c_m(\omega_m)}{12}\deleted{\unsure{\in\mathscr T_X^{*\delta^2}}}=\frac{\delta^2c_m(\omega_m)}{12}\in\mathscr T_X^{*\delta^2}\\
 \omega'_n&\unsure{\mapstochar\relbar|\longmapsto}\beta^2c'_n(\omega'_n)=\beta^2n^2\frac{c_n(\omega'_n)}{12}=\delta^2\frac{c_n(\omega'_n)}{12}\in\Gamma\mathscr T_X^{*\delta^2}
 \end{aligned} \]

\[ G'_{m,n}(\omega_m,\omega'_n)=\alpha^2c'_m(\omega_m)-\beta^2c'_n(\omega'_n)\in\Gamma\mathscr T_X^{*\delta^2} \]

\[ \left\{\begin{aligned}
 G_{\unsure{a|m}}(t,as)&=G(t,s)&a&\in\Gamma(S,\underline O_S^{*})\\
 G(t,s+b)&=\unsure{\bar G|G}(t,s)&b&\in\Gamma(S,\underline O_S)\\
 G(t,\tfrac1s)&=G(t,s)&s&\in\underline{\mathrm{Ét}}_S(X,G_{m,S})
 \end{aligned}\right. \]

% lower calculations written upside down
\[ \begin{aligned}
 f&\longmapsto af+b\\ f&\longmapsto af\\
 df^{\unsure{2|n}}&\qquad\left(\frac{df}{f}\right)^2\\
 \frac{df}{f}&=c\,\frac{dg}{g}\\
 \log f&=\unsure{c\log g^{c}|\log g^{c}}\,,\\
 f&=g^c\,
 \end{aligned} \]

\[ \begin{aligned}
 df=dg&\Rightarrow G(f,g)=0\\
 \left(\frac{df}{f}\right)^2=\left(\frac{dg}{g}\right)^{\unsure{2|1}}&\Rightarrow G(f,g)=0\\
 (df-dg)&\left(\left(\frac{df}{f}\right)^2-\left(\frac{dg}{g}\right)^2\right)
 \end{aligned} \]

% archive sheet number: 79
\scanpage{81}

\diagram{\begin{tikzcd}[ampersand replacement=\&,column sep=small,row sep=small]
D_X \& U(\unsure{x|n})\arrow[l,"{\mathbb R\text{-an}}"]\arrow[r,hook,"{\mathbb C\text{-an}}"]\arrow[d]\&
 B(x)\arrow[d,swap,"{\mathbb R\text{-an}}"] \&
 V(x)\arrow[l,hook',"{\mathbb C\text{-an}}"']\arrow[r,"{\mathbb R\text{-anal.}}"] \&
 \deleted{\unsure{D'_X|D'_x}}\arrow[d]\\
\&X_X\&\Sigma(x)\arrow[d,swap,"?"]\arrow[d,bend left,"{\mathbb R\text{-an}}"]\&\&\bar X_X\\
\&\&\widehat T_X\&\&
\end{tikzcd}}

\[ \begin{aligned}D_{\unsure{n|x}}&\subset T_X,\ \deleted{\unsure{n|x}}\ \supset D'_{\unsure{a|n|x}}\\ &\qquad\unsure{\shortparallel|?|\simeq}\\ &\qquad\bar D_n.\end{aligned} \]

\deleted{\[ \begin{aligned}
 \widetilde X_{(x)}\times_{(X,\Pi(x))}\overline X_{(x)}&=\qquad\deleted{\unsure{X}}\\
 X(x)\qquad\pi_1(X,x)&\quad\text{\unsure{n$^\circ$ d'elts|n$^\circ$ d'ents|ses éléments}}\\
 \widetilde X&\\ \Pi_x\times\bar\Pi_x&\qquad\unsure{.}\\
 \widehat\tau_X^*+O_X&\\
 \unsure{\mathcal V|\mathcal J}_X+\underline O_X&=\Omega^{-1}\oplus\Omega^0
 \end{aligned} \]}

\diagram{\begin{tikzcd}[ampersand replacement=\&,column sep=small,row sep=small]
 \bar X_X\&\Sigma(x)\&\bar X_X\\
 U(x)\arrow[u,"{\mathbb C\text{-an}}"]\arrow[d,swap,"{\mathbb R\text{-anal.}\ {?}}"]\&
 B(x)\arrow[l]\arrow[r]\arrow[u]\arrow[d,"{?\ \mathbb R\text{-anal.}}"]\&
 V(x)\arrow[u]\arrow[d,"{?\ \mathbb R\text{-anal.}}"]\\
 D_X\arrow[d]\&\widehat T_X\arrow[l]\arrow[r]\arrow[d]\&D''_X\\
 \deleted{\unsure{?}}\&\Sigma_X\&\\
 \&=\&
 \end{tikzcd}}

$\mathrm{GP}(1,\mathbb C)$

\deleted{$B(x)_y=$ \unsure{classe|classes} d'\unsure{équiv.|équivo.} \unsure{de couples|du couple} $(\unsure{\ell|c},\xi)$}

\deleted{\[ \begin{aligned}
 \xi\in\widehat T_{X,\unsure{a|x}} &\quad \ell\in\Pi(x,y)\\
 \deleted{(\ell,\xi}\ \bigl((c\lambda,\xi)&\simeq(c,\lambda.\xi)\\
 &(\lambda\in\Pi(x))
 \end{aligned} \]}

\[ \begin{aligned}
 \check E\otimes E&=(\Omega^1+\Omega^0)\otimes(\Omega^{-1}\oplus\Omega^0)\\
 &=\begin{pmatrix}\Omega^0&\Omega^{-1}\\ \Omega^1&\Omega^0\end{pmatrix}
 \end{aligned} \]

\[ \begin{aligned}
 \deleted{\underbrace{\check E\otimes E}\otimes\Omega^1\simeq\begin{pmatrix}\Omega^1&\Omega^0\\\Omega^2&\Omega^1\end{pmatrix}}\\
 \mathcal V\otimes\Omega^1\simeq\begin{pmatrix}\Omega^1&\Omega^0\\\Omega^2&-\Omega^1\end{pmatrix}
 \end{aligned} \]

\[ \deleted{\begin{pmatrix}1&0\\\unsure{2g-2}&1\end{pmatrix}/\unsure{2g}} \]

\[ \begin{aligned}
 f,g\qquad & \alpha^2=1\quad \alpha=\pm1 \quad\left\{\begin{array}{l}g+1\\g\end{array}\right.\qquad 3g-3\qquad 1\ \big/\ 4g-2\\
 &\qquad\qquad\qquad\qquad \deleted{\unsure{2|2g}}+g
 \end{aligned} \]

\[ \begin{aligned}
 f+\tfrac1f={}&\qquad\Omega^1=\omega^{\otimes2}\qquad =-f\ \unsure{\mathrm{défini}|df^{-1}}\\
 &\frac{f^2+1}{f}\quad\text{\uline{sans zéros}}\qquad\deleted{\unsure{d|\mathrm{et}}}\quad\frac{g}{g'}=\varphi\\
 f\longmapsto\tfrac1f\ \big/\quad & f+\tfrac1f=\frac{f^2+1}{f}=h\quad\unsure{?}\ \text{sans zéro.}\\
 & f^2-fh+1=
 \end{aligned} \]

\[ \begin{aligned} \frac{g'}g&=\frac1\varphi\\ \log g&=\int\frac1\varphi+\\ g&=\exp\int\frac1\varphi \end{aligned} \]

% archive sheet number: 80
\scanpage{82}

écart additif

\[ A_1(f,g)=dg/df=g'/f'\qquad\deleted{\unsure{A(f,g)}} \]

\[ \unsure{A|B}_n(f,g)=\left(\frac{dg}{df}\right)^{\unsure{n|m}}=\left(\frac{g'}{f'}\right)^n \]

écart affine

\[ \mathrm{Af}(f,g)=D_{dg}-D_{df}
 =\frac{d\,\frac{dg}{df}}{dg/df}=\frac{d\,\frac{g'}{f'}}{g'/f'}
 =\frac{dg'}{g'}-\frac{df'}{f'}=\left(\frac{g''}{g'}-\frac{f''}{f'}\right)dt \]

\diagram{\resizebox{\linewidth}{!}{\begin{tikzcd}[ampersand replacement=\&,column sep=small,row sep=small]
 \&\&\&\substack{\unsure{\mathscr E|\mathcal E}_X/(G_m\cdot G_a)(\mathbb C)\xrightarrow[c_{\mathrm{af}}(f)=D_{df}]{\sim}\mathrm{Conn}(\underline\omega_X^1)\\\text{torseur sous }\underline\omega_X^1}
 \arrow[d,"{D\mapsto c_D}"]\arrow[d,bend right,"?"']\\
 \&\&\substack{\mathscr E_X/(\mu_n\cdot G_a)(\mathbb C)\xrightarrow[c_{a,n}(f)=(df)^n]{\sim}\underline\omega_X^n\\\text{torseur sous }\underline O_X^*}
 \arrow[ur]\arrow[ur,bend right,swap,"{(n\ne0)\quad\psi\mapsto-\frac1n\frac{d\psi}{\psi}}"]\&
 \substack{\unsure{\mathscr P|\mathcal P|\mathcal E}_X/\mathrm{GP}(1,\mathbb C)\xrightarrow[c_{\mathrm{pr}}(f)=c_{D_{df}}]{\sim}\mathrm{Conn\,proj}(X)\\
 \underline\omega_X^2\quad(\text{torseur sous})}\\
 \&\substack{\mathscr E_X/G_a(\mathbb C)\xrightarrow[c_a(f)=df]{\sim}\underline\omega_X^{1*}\\\text{torseur sous }\underline O_X^*}
 \arrow[ur]\arrow[ur,bend right,swap,"{\varphi\mapsto\varphi^n}"]\arrow[dd,"?"]\arrow[dd,bend left,"{\mathrm{id.}}"]\&
 \&\boxed{c_{D+\omega_1}=c_D+2D\omega_1+\omega_1^2}\\
 \mathscr E_X\arrow[ur]\arrow[d,bend left,"{\exp\ \deleted{\unsure{\mathscr E_X}}}"]\arrow[d,dashed,swap,"{\deleted{\unsure{\mathscr E_X^*}}}"]\&\&
 \substack{\mathscr E_X^{\times}/(\mathbb Z/2\mathbb Z\cdot G_m)(\mathbb C)\simeq\underline\omega_X^{2*}\\
 \text{torseur sous }\underline O_X^*}\&\\
 \mathscr E_X^*\arrow[r]\&\substack{\mathscr E_X^*/G_m(\unsure{\mathbb C|m})\xrightarrow[c_m(f)=df/f]{\sim}\underline\omega_X^{1*}\\\text{torseur sous }\underline O_X^*}
 \arrow[ur,"{\unsure{\mathscr E_X}}"]\arrow[ur,bend right,swap,"{\varphi\mapsto\varphi^2}"]\&\&
 \end{tikzcd}}}

\[ \exp(\pm f+b)=\underbrace{\exp b}_{c}\,(\exp f)^{\pm1} \]

\[ \mathrm{Ecad}(f,g)=\frac{dg}{df}=\frac{g'}{f'}\in\Gamma\,\underline O_X^* \]

\[ \mathrm{Eca}\unsure{\hat\mu_n|\mu_n|f_n}(f,g)=\left(\frac{dg}{df}\right)^n=\left(\frac{g'}{f'}\right)^n\in\Gamma(\underline O_X^{\times}) \]

\[ \mathrm{Ecaf}(f,g)=\frac{d(dg/df)}{dg/df}=\frac{d(g'/f')}{g'/f'}=\frac{dg'}{g'}-\frac{df'}{f'}=\left(\frac{g''}{g'}-\frac{f''}{f'}\right)dt\in\Gamma(\underline\omega_X^1) \]

\[ \mathrm{Ecmu}\unsure{.|_1}(f,g)=\frac{dg/g}{df/f}=\frac{g'/g}{f'/f}\in\Gamma(\underline O_X^{\times}) \]

\[ \mathrm{Ecmu}_2(f,g)=\left(\frac{dg/g}{df/f}\right)^2=\left(\frac{g'/g}{f'/f}\right)^2\in\Gamma(\underline O_X^{\times}) \]

\[ \begin{aligned}\mathrm{Ecpr}(f,g)&=-2\frac{d^3g/df^3}{dg/df}+3\left(\frac{d^2g/df^2}{dg/df}\right)^2\\
 &=\left[\left(-2\frac{g'''}{g^{\prime}}+3\left(\frac{g''}{g'}\right)^2\right)-\left(-2\frac{f'''}{f'}+3\left(\frac{f''}{f'}\right)^2\right)\right]dt^2\in\Gamma(\underline\omega_X^2)
 \end{aligned} \]

\diagram{\resizebox{\linewidth}{!}{\begin{tikzcd}[ampersand replacement=\&,column sep=small,row sep=small]
 \mathscr E_X\arrow[r,"d"]\arrow[dddd,swap,"{\exp}"]\&
 {}[\mathscr E_X/G_a(\mathbb C)\xrightarrow[c_a(f)=df]{\sim}\underline\omega_X^{1*}]\arrow[r,"{\alpha_n}"]\arrow[dddd,swap,"?"]\arrow[dddd,bend left,"{\mathrm{id}}"]\&
 {}[\mathscr E_X/(\mu_n\cdot G_a)(\mathbb C)\xrightarrow[c_{a,n}(f)=(df)^n]{\sim}\underline\omega_X^{n*}]
 \arrow[r,"{\beta_n}"]\arrow[dddd,"{\mathrm{id}\;(n=2)}"]\&
 {}[\mathscr E_X/(G_m\cdot G_a)(\mathbb C)\xrightarrow[c_{\mathrm{af}}(f)=D_{df}]{\sim}\mathrm{Conn}(\underline\omega_X^1)]\arrow[dd,"{\gamma}"]\\
 \&\underline O_X^*\arrow[r,"{\varphi\mapsto\varphi^n}"]\arrow[dd,swap,"{\mathrm{id}}"]\&
 \deleted{\unsure{O}}\underline O_X^{\times}\deleted{\unsure{*}}\arrow[r,"{\psi\mapsto-\frac1n\frac{d\psi}{\psi}}"]\arrow[dd,swap,"{\mathrm{id}}"]\&\underline\omega_X^1\\
 \&\&\&{}[\mathscr P_X/\mathrm{GP}(1)(\mathbb C)\simeq\unsure{\mathrm{Conn}|\mathrm{Ann}}(X)]\\
 \&\underline O_X^{\times}\arrow[r,"{\varphi\mapsto\varphi^2}"]\&\underline O_X^{\times}\arrow[r,bend right]\&\underline\omega_X^2\\
 \mathscr E_X^*\arrow[r]\&{}[\mathscr E_X^{\times}/G_m(\mathbb C)\underset{c_m(f)=\frac{df}{f}}{\simeq}\underline\omega_X^{1*}]
 \arrow[r,"{\alpha'_2}"]\&{}[\mathscr E_X^*/(\mathbb Z/2\mathbb Z\cdot G_m)(\mathbb C)=\underline\omega_X^{2*}]
 \arrow[ur,bend right,swap,"{\delta}"]\&
 \end{tikzcd}}}

\[ \left\{\begin{aligned}
 \delta(\dot\omega_2)&=\gamma\beta_2(\omega_2)+\omega_2\\
 &=c_{\mathrm{pr}}(t)+\left[\left(-\deleted{\unsure{2}}\frac{f''}{f}+\frac{5}{\unsure{1|4}}\left(\frac{f'}{f}\right)^2\right)\right]dt^2\\
 \delta\alpha'_2(\omega_1)&=\gamma\beta_1(\omega_1)\deleted{\unsure{+}}+\omega_1^2\\
 &=c_{\mathrm{pr}}(t)+\left(f^2-2\frac{f''}{f}+3\left(\frac{f'}f\right)^2\right)dt^2
 \end{aligned}\right. \]

\[ \begin{aligned}
 \alpha_n(\omega_1)&=\omega_1^n\\
 \beta_n(\underset{\substack{\Vert\\f(dt)^n}}{\omega_n})&=\beta_1(\omega_n^{1/n})=\beta_1(f^{1/n}dt)=D_{dt}\deleted{\unsure{f}}+\frac1n\frac{df}{f}=D_{dt}+\frac1n\frac{f'}f\,dt
 \end{aligned} \]

\[ \gamma(D_{dt}+\underset{\substack{\Vert\\fdt}}{\omega_1})
 =\underset{\substack{\Vert\\c_{\mathrm{pr}}(t)}}{\gamma(D_{dt})}+2D_{dt}\omega_1+\omega_1^2
 =c_{\mathrm{pr}}(t)+2\frac{df}{\deleted{dt}}dt+f^2dt^2
 =c_{\mathrm{pr}}(t)+(2f'+f^2)dt^2 \]

\[ \begin{aligned}
 \gamma\beta_n(\omega_n)&=c_{\mathrm{pr}}(t)+\left(2\left(\tfrac1n\tfrac{f'}f\right)'+\left(\tfrac1n\tfrac{f'}f\right)^2\right)dt^2\\
 &=c_{\mathrm{pr}}(t)+\frac1{n^2}\left(-2n\frac{f''}{f}+\deleted{\unsure{n}}(2n+1)\left(\frac{f'}f\right)^2\right)dt^2\\
 \gamma\beta_1(\omega_1)&\qquad=c_{\mathrm{pr}}(t)+\left(-2\frac{f''}{f}+3\left(\frac{f'}f\right)^2\right)dt^2\\
 \gamma\beta_1d(f)&=c_{\mathrm{pr}}(t)+\deleted{\unsure{-2}}\left(-2\frac{f'''}{f'}+3\left(\frac{f''}{f'}\right)^2\right)dt^2
 \end{aligned} \]

% archive sheet number: 81
\scanpage{83}

\[ \unsure{X|Y}=\Omega Y \]

\diagram{\begin{tikzcd}[ampersand replacement=\&,column sep=small,row sep=small]
 k_1\&{}\in H^3(\pi_1(X),\pi_2(X))\overset{?}{=}0\\
 \exists\ k'_1\arrow[u,dashed]\&{}\in H^4(K(\pi_1(X),2),\pi_2(X))\arrow[u]\\
 k''_1\& H^5(K(\pi_1(X),3),\pi_2(X))\arrow[u]\\
 \&\unsure{\shortparallel|\wr\wr|\simeq}\\
 \&H^{5+\unsure{r|\gamma|i}}(K(\pi_1(X),3+r,\pi_2(X))
 \end{tikzcd}}

\diagram{\begin{tikzcd}[ampersand replacement=\&,column sep=small,row sep=small]
 \&Y'\arrow[d]\&\\
 \&K(\pi_1(X),2)\arrow[r,"{k'_1}"]\&K(\pi_2(X),4)\\
 S^1\wedge{}K(\pi_1(X),1)\arrow[r]\&
 K(\pi_1(X,2)\arrow[u,phantom,"{\unsure{\simeq|\sim}}" description]\&
 \end{tikzcd}}

\diagram{\begin{tikzcd}[ampersand replacement=\&,column sep=small,row sep=small]
 X=\Omega Y\arrow[d]\&\\
 K(\pi_1(X),1)\arrow[r,"{k_1}"]\& K(\pi_2(X),3
 \end{tikzcd}}

% archive sheet number: 82
\scanpage{84}

donc

\[ \left(\frac{df}{f}\right)^2=\frac{(df)^2}{f^2}=\left(\frac{{f'}^{\deleted{2}}}{f^{\deleted{2}}}\right)^2dt^2 \]

On a a) b) \deleted{\unsure{d|a}}\unsure{c)} d) e) -- On prendra

\[ E(f,g)=\left(\frac{dg}{g}\right)^2\mathbin{/}\left(\frac{df}{f}\right)^2=\frac{\left(\frac{dg}{df}\right)^2}{(g/f)^2} \]

\deleted{$\unsure{\mathscr E|\mathcal E}_X/G_a(\mathbb C)\unsure{\simeq|\in}$}

% archive sheet number: 83
\scanpage{86}

$\unsure{X|S|\text{\S}}$ schéma.

1.1. Données équivalentes

a) fibré en droites projectives \inserted{$P$} sur $X$ (fibré de Br. Sev. de rang 1) $f:P\longrightarrow X$

b) ``\deleted{\unsure{fibré de droites}} loc. libre \inserted{\unsure{rang|rg} 2 $E$} mod \deleted{\unsure{ten}}\unsure{tensoris.|torsion|tensorisation} par \unsure{module|modules} \unsure{inversible|inversibles}'' sur $X$

c) Algèbre d'Azumaya \inserted{$M$} de rang 4

\margin{si 2 inv. sur $X$}

d) Algèbre de Lie $\unsure{\mathfrak g|\mathcal G}$ forme \inserted{(fppf)} de $\mathfrak{gp}(1)_X$\\
 d') \rule{2.1cm}{0.3pt} $\mathfrak g'$ \rule{2.1cm}{0.3pt} \deleted{$\unsure{\mathfrak{pg}}$} $\unsure{\mathfrak{sgl}|\mathfrak{sl}}(1)_X$

e) forme $G$ de $\underline{\mathrm{GP}}(1)_X$\\
 \deleted{\unsure{e)}} \unsure{e$_1$)|e')} \rule{1.5cm}{0.3pt} $G'$ de $\underline{\mathrm{SL}}(2)_X$.

Si on a $P$, on en déduit $\bullet$ $E=f_*(L)$, $\underline L$ module inv. sur $P$ de degré 1 sur \unsure{ch.que|ch.} fibre ($\underline L$ défini loc. fpqc \inserted{sur $X$}, loc$^{\mathrm{lt}}$, mod tensorisation par $f^*(\lambda)$, $\lambda$ inv sur $X$)

\[ \begin{aligned}
 M&=\check E\otimes E\simeq\underline{\mathrm{End}}(E)\simeq E\otimes E\otimes(\overset2\Lambda E)^{-1}\quad\text{ext. de}\\
 &\underline O_X\;(\simeq\overset2\Lambda E\otimes(\overset2\Lambda E)^{-1})\quad\text{par }\Gamma^2E\otimes(\overset2\Lambda E)^{-1}
 \end{aligned} \]

\[ \text{ext. de }\underbrace{\mathrm{Sym}^2(E)\otimes(\overset2\Lambda E)^{-1}}_{f_*(L^{\otimes2})\otimes(\overset2\Lambda E)^{-1}\simeq f_*(\underbrace{L^{\otimes2}\otimes/\overset2\Lambda E}_{\underline t_{P/X}})}
 \text{ par }\underline O_X\;(\simeq\overset2\Lambda E\otimes(\overset2\Lambda E)^{-1}) \]

\margin{expliciter !}

\diagram{\begin{tikzcd}[ampersand replacement=\&,row sep=small]
 \mathrm{Sym}^2\check E\otimes\mathrm{Sym}^2 E\arrow[d]\\
 \mathrm{Sym}^2\check E\otimes\overset2\Lambda E\\
 \text{et}\\
 \Gamma^2\check E\otimes\Gamma^2 E\arrow[d]\\
 \Gamma^2\check E\otimes\overset2\Lambda E
 \end{tikzcd}}

\margin{\deleted{\unsure{des coefficients}} formes de Killing}

\diagram{\begin{tikzcd}[ampersand replacement=\&,column sep=small,row sep=tiny]
 \mathrm{Sym}^2\check E\otimes\mathrm{Sym}^2 E\arrow[r]\&(\overset2\Lambda\check E)^{\otimes2}\\
 \Gamma^2\check E\otimes\Gamma^2 E\arrow[ur]\&
 \end{tikzcd}}

\[ \begin{aligned}
 M&\overset{?}{\simeq}f_*(\mathcal P^1_{P/X}(\underline t_{P/X}))\quad\text{\unsure{additivement|additivem.|addition} ?}\\
 \mathfrak g&=M/\underline O_X \quad\text{avec crochet déduit du crochet }\inserted{xy-yx}\text{ dans }M\ [\underline{\mathrm{NB}}\ \underline O_X\text{ en est le centre}]\\
 &\simeq f_*(\underline t_{P/X})\quad [\text{avec crochet provenant du crochet dans }\underline t_{P/X}]
 \end{aligned} \]

\[ \begin{aligned}
 \mathfrak g'&\subset M=\unsure{E|\mathrm{tr}=0}\quad\text{sous-algèbre des éléments de trace réduite nulle}\\
 &(\mathfrak g'\longrightarrow\mathfrak g\ \text{can., iso si 2 inv. sur }X)\\
 \mathfrak g'&\subset\overset{\Vert}{\check E}\otimes E\nearrow\check E\quad\text{sous. algèbre des él. de trace nulle}
 \end{aligned} \]

\[ \begin{aligned}
 G&=\mathrm{Aut}_X P,\qquad G'=\widetilde G\quad(\text{rev. \unsure{univ.} de Chevalley})\\
 &&=\deleted{\unsure{\mathrm{Aut}}}\underline{\mathrm{SL}}(E)
 \end{aligned} \]

\margin{$\begin{aligned}G&\simeq\underline{\mathrm{Aut}}(E)/G_{mX}\\&\simeq\underline{\mathrm{Aut}}(M)\\&\simeq\underline{\mathrm{Aut}}(\mathfrak g)\\&\simeq\underline{\mathrm{Aut}}(\mathfrak g')\end{aligned}$}

$P=P(E)={}$ schéma des Idéaux à gauche qui sont ``spéciaux'', i.e. loc. fact. directs et ``non \uline{\unsure{triviaux, i.e aussi|triviaux, au sens}}'' \deleted{\unsure{non tr.}}

\[ \begin{aligned}
 &=\text{schéma des sous-Alg. de Borel de }\mathfrak g\\
 &=\rule{6cm}{0.3pt}\quad\mathfrak g'\\
 &=\rule{4cm}{0.3pt}\quad\text{Gr. de Borel de }G\\
 &=\rule{6cm}{0.3pt}\quad\text{--- }G'
 \end{aligned} \]

\margin{! attention à la Car. (notamment Car. 2 !)}

% archive sheet number: 85
\scanpage{87}

\[ \begin{aligned}
 E&=\text{module d'une ``repr. \unsure{noire|naïve|vraie}'' de }\deleted{\unsure{M}}M,\\
 &\hspace{4cm}\text{de }\unsure{\mathfrak g|\mathcal G}\\&\hspace{4cm}\text{de }\mathfrak g'\\
 &\hspace{4cm}\text{de }G\\&\hspace{4cm}\text{de }G'
 \end{aligned} \]

\[ \begin{aligned}\mathfrak g&=\underline{\mathrm{Lie}}\,G\\ \mathfrak g'&=\underline{\mathrm{Lie}}\,G'\end{aligned} \]

1.2. Données équivalentes

a) fibré \deleted{\unsure{projectif}}\inserted{en droites} proj. $P$ sur $X$, avec section marquée $\sigma$

b) Structure d'extension
\[ 0\longrightarrow\Omega\longrightarrow\mathscr L\longrightarrow\underline O_X\longrightarrow0 \]
$\Omega$ inversible

c) Algèbre d'Azumaya de rg 4 $M$, avec idéal ``spécial'' \deleted{\unsure{$\mathfrak a$}} \unsure{$\mathcal V_\ell$|$\mathcal V_L$|$\mathfrak a$}

d) $\mathfrak g$ avec ss.-Alg. de Borel $\mathscr L$\\
d') $\mathfrak g'$ \rule{4cm}{0.3pt} $\mathscr L'$\\
e) $G$ \rule{2cm}{0.3pt} Gr. \rule{1.5cm}{0.3pt} $B$\\
e') $G'$ \rule{2cm}{0.3pt} gr \rule{1.5cm}{0.3pt} $B'$

% slanted
\margin{\unsure{identifier|vérifier} ! \uline{NB} $\mathfrak g/\Omega\simeq\check{\mathscr L}$ (par Killing) $\simeq\mathscr L\otimes\det\mathscr L^{-1}\simeq\mathscr L\otimes\Omega^{-1}$ avec filtration indiquée}

\deleted{f) Une $\underline O_X$-Algèbre augmentée $B_2$ \deleted{\unsure{trinquée}} \unsure{d'algèbre|d'échelon} \deleted{\unsure{d'idéaux}} \inserted{tronquée à l'}2, de ``rang'' 1}

\inserted{à a) est associé $I=\underline L$} faisceau inv. de degré relatif 1 \uline{rigidifié} par $\sigma^*(\underline L)\simeq\underline O_X$. Alors

\[ \begin{aligned}
 \bullet\quad\mathscr L&=f_*(\underline L),\quad\text{avec structure d'ex. évidente}\\
 &\simeq f_*(\underline L/\mathcal I\underline L)\qquad [\mathcal I=\text{\uline{Idéal}}_{\,\unsure{\sigma}}\text{ de }\sigma(X)\text{ dans }P]\\
 &\simeq\sigma^*(P^1(\underline L))
 \end{aligned} \]

$P=P(\mathscr L)$ avec section provenant du quotient inv. $\underline O_X$ de $\mathscr L$ $={}$ \unsure{cptfion|chpfion} projective du torseur sous $\unsure{\mathbb V|\mathbb W}(\Omega)$ défini par l'extension \inserted{de $\mathscr L$} $={}$ $\unsure{f_{\mathrm{nd}}|\mathrm{End}}\,\mathscr L$, avec idéal augmentation de $\Omega$

% slanted
\margin{(\unsure{adjoint} \deleted{du} prolongé)}

$\mathfrak g=M/\underline O_X\simeq\mathrm{Sym}^2(\mathscr L)\otimes\Omega^{-1}$ ($\underline{\mathrm{NB}}$ $\det\mathscr L\simeq\Omega$) Or $\mathrm{Sym}^2\mathscr L$ est filtré par (décr.) $\underline O_X,\Omega,\Omega^2$, donc $\mathrm{Sym}^2\mathscr L\otimes\Omega^{-1}$ filtré par (croiss.)

\diagram{$\begin{array}{ccc}
 \underbrace{\Omega}&\underline O_X&\Omega^{-1}\\
 \Vert&\Vert&\Vert\\
 \mathrm{Fil}_1&\mathrm{Fil}_2/\mathrm{Fil}_1&\mathrm{Fil}_3/\mathrm{Fil}_2.\\
 \Vert&\Vert&\Vert\\
 \mathrm{gr}_1&\mathrm{gr}_2&\mathrm{gr}_3
 \end{array}$}

% slanted
\margin{expliquer aussi \unsure{comment}, partant de $(\mathfrak g,\mathscr L)$, on trouve $\Omega\hookrightarrow\mathscr L\longrightarrow\underline O_X\longrightarrow0$ $(\Omega=\mathscr L_{\unsure{u|n}},\ldots)$}

\[ \begin{aligned}
 \text{où}\quad \mathrm{Fil}_2&=\Omega^{-1}\otimes\underbrace{(\mathrm{Ker}\ \mathrm{Sym}^2\mathscr L\longrightarrow\mathrm{Sym}^2 O)}_{\simeq\Omega\otimes\mathscr L}\\
 &\simeq\mathscr L
 \end{aligned} \]

donc on a
\[ \underset{\Omega}{\underbrace{\Omega}}\subset\underset{\underline O_X}{\underbrace{\mathscr L}}\subset\underset{\Omega^{-1}}{\underbrace{\mathfrak g}} \]
et $\mathscr L$ est une \uline{ss-Alg. de Borel} ($\Omega$ sa partie unipotente)

\margin{\uline{NB} pour forme de Killing, $\Omega$ et $\mathscr L$ orth. l'un de l'autre, si \unsure{2 inv.}}

On n'explicite pas $(\mathfrak g',\mathscr L')$ \unsure{$\overset{2}{\simeq}$} $(\mathfrak g,\mathscr L)$ si 2 inv), pour $(G,B)/(G',B')$ on a $G'=\underline{\mathrm{Sl}}_{\Omega_S}\mathscr L$, $B'\inserted{\subset G'}={}$ stabilisateur de $\Omega$ $G,B$ déduit par passage au quotient par centre $\mu_2$

[ \uline{\unsure{1.3.|1.4.}} Supposons $X/S$, et considérons des \deleted{connexions sur}\inserted{rel. $S$} les \unsure{Objets|objets} du type envisagé dans 1.1, mais moyennant \deleted{\unsure{l'existen}}\inserted{la donnée de} la structure 1.2. Ainsi, on considère des couples $(P,\sigma,c)$

% slanted
\margin{\uline{plus bas}}

\[ \left\{\begin{aligned} &P\quad\text{fibré en droites projectives sur }X\\ &\sigma\quad\text{section}\\&c\quad\text{connexion de }P\text{ rel. à }S\end{aligned}\right. \]

(mais on ne suppose pas $\sigma$ plate) \unsure{Il nous|Il nous y} sera particulièrement commode d'expliciter en termes de $(\mathscr L,\Omega)$ i.e.
\[ 0\longrightarrow\Omega\longrightarrow\mathscr L\longrightarrow\underline O_X\longrightarrow0 \]
(qui donne $(P,\sigma)$, \inserted{où $(\mathfrak g,\mathscr L)$, comme expliqué} et d'une connexion de $\mathfrak g$. \deleted{On note qu'une} ]

% archive sheet number: 86
\scanpage{88}

\deleted{\unsure{Splitting|Expliquons|Explicitions|Splittings}}

\diagram{\resizebox{\linewidth}{!}{\begin{tikzcd}[ampersand replacement=\&,row sep=small]
 \substack{\text{\uline{1.3.}\quad fibré }\deleted{\text{proj}}\text{\unsure{de droites|droites|en droites} projectives sur }X\text{ avec}\\
 \text{\uline{deux} sections \uline{disjointes} }\sigma,\sigma'}\arrow[d,leftrightarrow]\\
 \substack{\text{fibré inversible }\Omega\quad(P=\mathbb P(\Omega+\underline O_X)\simeq\overline{\unsure{\mathbb W|\mathbb V}(\Omega)}\\
 \text{clôture projective})\\ \text{avec }\sigma=\text{section à l'infini}\\\sigma'=\text{section nulle}}\arrow[d,leftrightarrow]\\
 {}[\,M\ \substack{\text{somme directe de deux Idéaux}\\\text{spéciaux ---}}\,]\arrow[d,leftrightarrow]\\
 (\unsure{\mathfrak g|\mathcal G})\text{ avec }\deleted{\text{\unsure{deux sous, sous-alg.}}}\inserted{\mathscr L\text{ et }\unsure{\text{sections s de B}}},\
 \text{\unsure{s'encontrant|s'encontr.} sur 1 \unsure{ss. Cartan|SGA 3}}
 \end{tikzcd}}}

soit décomposition en somme

\[ \mathfrak g\simeq\underbrace{\underset{\substack{\Vert\\\mathscr L_u}}{\Omega}+
 \underset{\substack{\Vert\\\underline O_X}}{\underline O_X^0}}_{\mathscr L}+\Omega^{-1} \]

\[ \begin{aligned}
 [\omega_1,\lambda]&=-[\lambda,\omega_1]=-\lambda\omega_1\\
 {}[\omega_{-1},\lambda]&=-[\lambda,\omega_{-1}]=\lambda\omega_{-1}\\
 {}[\omega_1,\omega_{-1}]&=-[\omega_{-1},\omega_1]=2\omega_1\omega_{-1}
 \end{aligned} \]

% slanted
\margin{\deleted{$\unsure{\mathfrak g,\mathscr L;\mathfrak g,\mathscr L'\to\mathfrak g',\mathscr L'}$} sans hyp. de car. $\ne2\unsure{,3|!}$ \uline{\unsure{essentiellement}}}

soit $\mathfrak g$ avec $t\subset\mathfrak g$ sous-alg. de Cartan, et le choix d'une des deux sous-Algèbres \unsure{unipotentes|nilpotentes} $\Omega,\Omega'$ normalisées par \uline{$t$} (ce qui fixe l'iso $t\simeq\underline O_X\simeq\underline{\mathrm{Hom}}_{\mathcal O}(\Omega,\Omega)$)

$\updownarrow$ $\mathfrak g'$ avec données analogues (Mais dans la Tble de multiplication dans les derniers mbs il faut
\[ \begin{aligned}-2\lambda\omega_{\unsure{1}},&\\2\lambda\omega_{-1}&\\\omega_1\omega_{-1}&\end{aligned} \]
)

\diagram{\begin{tikzcd}[ampersand replacement=\&,row sep=small]
 {}\arrow[d,leftrightarrow]\\
 G\text{ avec couple de Killing}\arrow[d,leftrightarrow]\\
 G'\text{ avec couple de }\unsure{\text{Killing}|\text{Ki Ki}}\ \unsure{\triangleright|D}
 \end{tikzcd}}

% archive sheet number: 87
\scanpage{89}

\uline{1.4 suite}

Choisissons d'abord un splitting de \inserted{$\mathscr L$ i.e. de} $\unsure{\mathfrak g|\mathcal G}$, un

\[ \mathfrak g\simeq\Omega+\underset{\substack{\Vert\\\underline O_X}}{\deleted{\unsure{O}}\Omega^0}+\Omega^{-1} \]

d'où, posant $\omega=\underline\Omega^1_{X/S}$
\[ \mathfrak g\otimes\omega\simeq\Omega\otimes\omega+\omega+\Omega^{-1}\otimes\omega \]

\deleted{\unsure{On en a une}} Choisissons-en loc. une section $\deleted{\unsure{\omega}}\delta$ de $\Omega$ comme base. Une connexion (sans courbure) $c_0$ est obtenue par

\[ \begin{aligned}
 \mathfrak g&\xrightarrow{c_0}\mathfrak g\otimes\omega\\
 c_0(\delta)&=c_0(1)=c_0(\delta^{-1})=0
 \end{aligned} \]

i.e.
\[ c_0\bigl(\underbrace{\lambda\delta+\mu.1+\nu\delta^{-1}}_{\substack{\updownarrow\\(\lambda,\mu,\nu)}}\bigr)=\underset{\substack{\updownarrow\\(d\lambda,d\mu,d\nu)}}{\delta\,d\lambda+1\,d\mu+\delta^{-1}d\nu} \]

Une connexion générale s'obtient alors en \unsure{ajoutant|rajoutant} une section de
\[ \mathrm{Dér}(\mathfrak g,\mathfrak g)\otimes\omega\simeq\mathfrak g\otimes\omega \]
(\uline{NB} toute dérivation de $\mathfrak g$ est intérieure !)

Soit \inserted{$\varpi_{\unsure{*|\bullet}}={}$} $\varpi_2\oplus\varpi_1\oplus\varpi_0\in\Gamma\mathfrak g\otimes\omega$ i.e.
\[ \begin{aligned}
 \varpi_2&\in\Gamma(\Omega\otimes\omega)\\
 \varpi_1&\in\Gamma(\omega)\\
 \varpi_0&\in\Gamma(\Omega^{-1}\otimes\omega)
 \end{aligned} \]

on aura donc \unsure{pour|par} connexion $c^{\deleted{\unsure{\varpi}}}_{\varpi_*}$ définie par $c_0$ et $\varpi_*$ 
\[ c_{\varpi_*}\bigl(\underbrace{\lambda\delta+\mu.1+\nu\delta^{-1}}_{(\lambda,\mu,\nu)}\bigr)\overset{?}{=}\deleted{\unsure{\delta\,d\lambda+1\,d\mu+\delta^{-1}d\nu}} \]

\[ \unsure{+|\ast} \]

% archive sheet number: 88
\scanpage{90}

\[ c_{\varpi_*}(\underset{\substack{\Vert\\\omega_{\unsure{0|1}}}}{\lambda\delta})
 =\underbrace{\delta\,d\lambda}_{\text{poids 2}}+
 \deleted{\unsure{\lambda}}\bigl[\underset{\substack{\Vert\\0}}{[\varpi_2,\lambda\delta]}+
 \underset{\substack{\Vert\\\underbrace{+\varpi_1(\lambda\delta)}_{\text{poids 2}}}}{[\varpi_1,\lambda\delta]}+
 \underset{\substack{\phantom{\Vert}\\\deleted{\unsure{-}}\,2\varpi_0(\lambda\delta)}}{[\varpi_0,\lambda\delta]}\bigr] \]

\[ \begin{aligned}
 c_{\varpi_*}(\underset{\substack{\Vert\\\lambda\delta}}{\omega_1})&=\overbrace{\delta\,d\lambda}^{2}+\overbrace{\underset{\substack{\Vert\\\lambda\delta}}{\omega_1}(\overbrace{\varpi_1}^{2}-2\varpi_0)}^{1}\\
 c_{\varpi_*}(\underset{\substack{\Vert\\\mu.1}}{\omega_0})&=\overbrace{d\mu}^{1}+\overbrace{\omega_0(\overbrace{-\varpi_2}^{2}+\varpi_0)}^{0}\\
 c_{\varpi_*}(\underset{\substack{\Vert\\\nu\delta^{-1}}}{\omega_{-1}})&=\overbrace{\delta^{-1}d\nu}^{0}+\overbrace{\omega_{-1}(\overbrace{2\varpi_2}^{1}-\varpi_1)}^{0}
 \end{aligned} \]

\[ \begin{aligned}
 \unsure{\mathfrak g|\mathcal G}&:\quad \Omega,\ \underset{\substack{\Vert\\\underline O_X}}{\Omega^0},\ \Omega^{-1}\\
 \mathfrak g\otimes\omega&:\quad\Omega\otimes\underline\omega\quad\underline\omega\quad\Omega^{-1}\otimes\underline\omega
 \end{aligned} \]

On trouve \unsure{les conclusions suivantes|la conclusion suivante|la condition suivante} (\unsure{indépendantes|indépend$^{\mathrm{tes}}$|indépendante} du choix des splittings \ldots !)

\uline{tte connexion $c$ satisfait} $c(\mathrm{Fil}_i\,\mathfrak g)\subset\mathrm{Fil}_{i+1}(\mathfrak g\otimes\underline\omega)$ par passage au quotient, il induit
\[ c_{i,i+1}:\ \deleted{\unsure{\mathrm{gr}(\mathfrak g)}}\,\mathrm{gr}_i\,\mathfrak g\longrightarrow\mathrm{gr}_{i+1}(\mathfrak g\otimes\omega)\simeq\mathrm{gr}_{i+1}(\mathfrak g)\otimes\omega \]
qui est \uline{linéaire}, de façon précise

\[ \begin{aligned}
 c_{1,2}:\mathrm{gr}_1\longrightarrow\mathrm{gr}_2,\quad \Omega\longrightarrow\omega,&\quad\text{i.e. }c_{12}\in\Gamma\,\Omega^{-1}\otimes\underline\omega\\
 c_{2,3}:\mathrm{gr}_2\longrightarrow\mathrm{gr}_3\quad \underline O_X\longrightarrow\Omega^{-1}\otimes\omega,&\quad\text{i.e. }c_{23}\in\Gamma\,\Omega^{-1}\otimes\underline\omega
 \end{aligned} \]

on a [$c_{12}=-2\varpi_0$, $c_{2,3}=\varpi_0$, donc ]
\[ \boxed{c_{1,2}=-2c_{2,3}} \]

\uline{NB} On trouve que la connexion respecte la filtration ssi $\underline{c_{2,3}}=0$ (d'où $c_{12}=0$) [cela signifie $\varpi_0=0$] -- \deleted{\unsure{celles qui la}}\inserted{ce sont aussi} celles qui respectent le Borel $\mathscr L\subset\mathfrak g$, ou encore \deleted{telles} qui respectent la section $\sigma$ de $P$, i.e. pour lesquelles $\sigma$ est horizontale. Elles \unsure{sont|se} \deleted{\unsure{déterminées}}\unsure{diffèrent de|différenciées} une section de $\mathscr L\otimes\omega$. \deleted{\unsure{On obtient localement $f$, on chois}}

\margin{Re \uline{NB} On peut aussi considérer $c^\Omega:\underset{\substack{\Vert\\\check{\mathscr L}\simeq\mathscr L\otimes\Omega^{-1}}}{\mathfrak g/\Omega}\longrightarrow\Omega^{-1}\otimes\underline\omega$ induit par $c$, induisant évidemment sur $\Omega^0\simeq\underline O_X\subset\mathscr L\otimes\Omega^{-1}$, sa connaissance est \uline{équivalente} à celle des $c_\Omega$ et 2 signes, ets dans le cas général (\unsure{peut-être en|précisément} car. 2 ? C'est pas \unsure{sur} plus précis en}

\uline{NB} Les connexions sur $\mathfrak g$ compatibles avec crochet \deleted{\unsure{(plus}} (si 2 inv. -- sinon il faut parler des connexions sur $P$, ou $G$ etc) forment un torseur sous $\mathfrak g\otimes\underline\omega$. Or
\[ \mathfrak g\otimes\underline\omega\longrightarrow \Omega^{-1}\otimes\underline\omega \]
d'où un torseur sous $\Omega^{-1}\otimes\underline\omega$. \uline{Ce dernier} est splitté canoniq$^{\mathrm{t}}$, i.e. on a
\[ \underline{\mathrm{Conn}}\xrightarrow{\varpi_0}\Omega^{-1}\otimes\underline\omega \]
compatible avec \unsure{ce|$\varpi_0$}, et si $c$ section de $\underline{\mathrm{Conn}}$, $\varpi_0(c)$ détermine les hom $\mathrm{gr}_i(\mathfrak g)\to\mathrm{gr}_{i+1}(\mathfrak g\otimes\underline\omega)$

\[ \begin{aligned}
 \Omega=\mathrm{gr}_1(\mathfrak g)&\xrightarrow[-2\varpi_0]{}\omega=\mathrm{gr}_2(\mathfrak g\otimes\Omega^{\unsure{1}})\\
 \Omega^0=\underline O_X=\mathrm{gr}_2(\mathfrak g)&\xrightarrow[\varpi_0]{}\Omega^{-1}\otimes\underline\omega=\mathrm{gr}_3(\mathfrak g\otimes\underline\omega)
 \end{aligned} \]

Si on considère les $c$ pour lesquels $\varpi_0(c)\in\Gamma\underset{\substack{\Vert\\\mathrm{Hom}(\Omega,\underline\omega)}}{\Omega^{-1}\otimes\underline\omega}$ est donné \inserted{$i$}, on donne une restriction du Gr. struct. de $\underline{\mathrm{Conn}}$ à $\mathscr L\otimes\underline\omega$, d'où torseur $\mathrm{Conn}_i$ sous $\mathscr L\otimes\underline\omega$. Or on a
\[ \mathscr L\otimes\underline\omega\longrightarrow\underline\omega \]
d'où un torseur associé sous $\underline\omega$ : les conn. sur $\mathfrak g$ compatibles avec $i$, \uline{modulo} \inserted{op. d'}ajouter une section de $\Omega\otimes\underline\omega\subset\mathscr L\otimes\underline\omega\subset\mathfrak g\otimes\underline\omega$. C'est aussi (\deleted{\unsure{raisonnant}}) le torseur des restrictions à $\Omega$ des conn. compatibles avec $i$, en ce cas des hom additifs op. diff. d'ordre \uline{1} (rel $S$

\[ \Omega\xrightarrow{c_\Omega}\mathscr L\otimes\omega \]
telles que la composée \deleted{$\unsure{\Omega\to\Omega\otimes\omega}$} \inserted{$\rightarrow\mathscr L\otimes\omega\longrightarrow\omega$} soit linéaire (et égale à $-2i$), et que l'on ait
\unsure{sc|i.e.|($\ast\ast$)}
\[ c_\Omega(fw)-fc_\Omega(w)=w\otimes df\qquad \begin{aligned}w&\in\Gamma\Omega\\f&\in\Gamma\underline O_X\end{aligned} \]

\margin{TSVP}

% archive sheet number: 89
\scanpage{91}

\deleted{(\unsure{indp. de $i$})}

La donnée d'un tel $c_\Omega$ équivaut à une restriction du Groupe structural de $\underline{\mathrm{Conn}}_i$ \deleted{\unsure{$\mathscr L$|$\varepsilon$}} à $\deleted{\unsure{\omega}}\Omega\otimes\underline\omega$, ou encore (\unsure{indépendamment|indépendant} de $i$) la donnée d'un $c_\Omega$ pour lequel ($\ast$) est lin et ($\ast\ast$) est \unsure{vér.|vérifiée|vraie}, équivaut à une restriction du Groupe structural $\unsure{\mathfrak g|\mathcal G}\otimes\omega$ de $\underline{\mathrm{Conn}}$ à $\Omega\otimes\underline\omega$..

On s'intéresse au cas où \deleted{\unsure{$c_i$|$\underline{\mathrm{Conn}}_i$}} $c_\Omega$ est tel que la composée
\[ (-2i):\Omega\longrightarrow\mathscr L\otimes\underline\omega\unsure{\relbar|\longrightarrow}\underline\omega \]
soit un \uline{iso}. \inserted{(\uline{ce qui implique 2 inv.})} Alors \deleted{\unsure{le}} le triplet

$(\mathfrak g,\mathscr L,c_\Omega)$ (\deleted{ou} ce qui revient au $=$, $(\mathscr L\supset\Omega,c_\Omega)$) est défini à isom unique près par $X/S$ (et existe ssi $\underline\omega=\underline\Omega^1_{X/S}$ est inversible) comme
\[ \mathscr L\simeq\mathcal P^1_{X/S}(\underline\omega)\otimes\underline\omega^{-1}\quad(\text{donc }\mathscr L\otimes\underline\omega=\mathcal P^1_{X/S}(\underline\omega)) \]

et $c_\Omega:\Omega\simeq\underline\omega\longrightarrow\mathscr L\otimes\underline\omega\simeq\mathcal P^1_{X/S}(\underline\omega)$ n'est autre que $d^1_{X/S,\deleted{\unsure{\omega}}\underline\omega}$ --

% archive sheet number: 90
